\input{_preamble_ftn.tex}
\begin{document}
\hd{FTN P1 — tematski listovi (rešenja)}{10 slotova · potpuna rešenja sa rezultatima}

\section*{Slot 1. Kompleksni brojevi — rešenja}
\begin{enumerate}[label=\arabic*.]
\item \textbf{a)} $|z|=\sqrt{(-3)^2+(3\sqrt{3})^2}=\sqrt{9+27}=6.$ Broj se nalazi u drugom kvadrantu ($\operatorname{Re}z<0,\ \operatorname{Im}z>0$), $\tan\varphi=\dfrac{3\sqrt{3}}{-3}=-\sqrt{3}$, pa je $\arg z=\pi-\dfrac{\pi}{3}=\dfrac{2\pi}{3}.$ Rezultat: $z=6\left(\cos\dfrac{2\pi}{3}+i\sin\dfrac{2\pi}{3}\right).$

\textbf{b)} Za $1+i$: moduo $\sqrt{2}$, argument $\dfrac{\pi}{4}$. Za $\sqrt{3}-i$: moduo $2$, argument $-\dfrac{\pi}{6}$. Pri deljenju se moduli dele, a argumenti oduzimaju:
\[ |z|=\frac{\sqrt{2}}{2},\qquad \arg z=\frac{\pi}{4}-\left(-\frac{\pi}{6}\right)=\frac{5\pi}{12}\in(-\pi,\pi]. \]
Rezultat: $z=\dfrac{\sqrt{2}}{2}\left(\cos\dfrac{5\pi}{12}+i\sin\dfrac{5\pi}{12}\right).$

\textit{Rešenje:} a) $|z|=6,\ \arg z=\dfrac{2\pi}{3},\ z=6\left(\cos\dfrac{2\pi}{3}+i\sin\dfrac{2\pi}{3}\right)$;\ b) $|z|=\dfrac{\sqrt{2}}{2},\ \arg z=\dfrac{5\pi}{12},\ z=\dfrac{\sqrt{2}}{2}\left(\cos\dfrac{5\pi}{12}+i\sin\dfrac{5\pi}{12}\right).$

\item \textbf{a)} $\dfrac{\sqrt{3}+i}{2}=\cos\dfrac{\pi}{6}+i\sin\dfrac{\pi}{6}$ (moduo $1$). Po Moavru ugao $2026\cdot\dfrac{\pi}{6}=\dfrac{1013\pi}{3}.$ Redukcija: $\dfrac{1013\pi}{3}=2\pi\cdot168+\dfrac{5\pi}{3}.$ Znači rezultat je jednak $\cos\dfrac{5\pi}{3}+i\sin\dfrac{5\pi}{3}=\dfrac{1}{2}-\dfrac{\sqrt{3}}{2}\,i.$

\textbf{b)} $1-i=\sqrt{2}\left(\cos\left(-\dfrac{\pi}{4}\right)+i\sin\left(-\dfrac{\pi}{4}\right)\right).$ Tada je moduo $(\sqrt{2})^{2026}=2^{1013}$, ugao $2026\cdot\left(-\dfrac{\pi}{4}\right)=-\dfrac{1013\pi}{2}\equiv-\dfrac{\pi}{2}\ (\mathrm{mod}\ 2\pi).$ Zato je $(1-i)^{2026}=2^{1013}\left(\cos\left(-\dfrac{\pi}{2}\right)+i\sin\left(-\dfrac{\pi}{2}\right)\right)=-2^{1013}\,i.$

\textbf{v)} $w=-8+8\sqrt{3}\,i$: moduo $\sqrt{64+192}=16$, argument $\dfrac{2\pi}{3}$. Koreni imaju moduo $\sqrt[4]{16}=2$ i argumente $\dfrac{1}{4}\left(\dfrac{2\pi}{3}+2\pi k\right)=\dfrac{\pi}{6}+\dfrac{\pi k}{2},\ k=0,1,2,3$, to jest uglove $\dfrac{\pi}{6},\ \dfrac{2\pi}{3},\ \dfrac{7\pi}{6},\ \dfrac{5\pi}{3}$:
\[ z_0=\sqrt{3}+i,\quad z_1=-1+\sqrt{3}\,i,\quad z_2=-\sqrt{3}-i,\quad z_3=1-\sqrt{3}\,i. \]

\textit{Rešenje:} a) $\dfrac{1}{2}-\dfrac{\sqrt{3}}{2}\,i$;\ b) $-2^{1013}\,i$;\ v) $\sqrt{3}+i,\ -1+\sqrt{3}\,i,\ -\sqrt{3}-i,\ 1-\sqrt{3}\,i.$

\item \textbf{a)} $z\bar z=|z|^2=3^2+(-2)^2=13.$ Dalje $\dfrac{\bar z}{z}=\dfrac{3+2i}{3-2i}=\dfrac{(3+2i)^2}{13}=\dfrac{5+12i}{13},$ i $\dfrac{1}{z}=\dfrac{\bar z}{|z|^2}=\dfrac{3+2i}{13}.$

\textbf{b)} $\dfrac{4+7i}{2+i}=\dfrac{(4+7i)(2-i)}{5}=\dfrac{15+10i}{5}=3+2i;$ $\dfrac{1-3i}{1+i}=\dfrac{(1-3i)(1-i)}{2}=\dfrac{-2-4i}{2}=-1-2i.$ Zbir: $w=(3+2i)+(-1-2i)=2.$

\textit{Rešenje:} a) $z\bar z=13,\ \dfrac{\bar z}{z}=\dfrac{5+12i}{13},\ \dfrac{1}{z}=\dfrac{3+2i}{13}$;\ b) $\operatorname{Re}w=2,\ \operatorname{Im}w=0.$

\item \textbf{a)} Diskriminanta $D=(3+i)^2-4(2+2i)=9+6i-1-8-8i=-2i.$ Nađimo $\sqrt{-2i}$: za $(x+yi)^2=-2i$ imamo $x^2-y^2=0,\ 2xy=-2$, odakle $x=1,\ y=-1$, to jest $\sqrt{D}=\pm(1-i).$ Tada je $z=\dfrac{(3+i)\pm(1-i)}{2}$:
\[ z_1=\dfrac{4}{2}=2,\qquad z_2=\dfrac{2+2i}{2}=1+i. \]
Provera (Vijetove formule): $z_1+z_2=3+i,\ z_1z_2=2+2i.$

\textbf{b)} Zamenimo $z=1+2i$: $z^2=-3+4i$, $(a-i)(1+2i)=(a+2)+(2a-1)i$. Zbir sa $b+3i$ jednak je nuli:
\[ a+b-1=0,\qquad 2a+6=0. \]
Odatle $a=-3$, $b=1-a=4.$

\textit{Rešenje:} a) $z_1=2,\ z_2=1+i$;\ b) $a=-3,\ b=4.$

\item Napišimo kvadrat modula:
\[ |z|^2=(t+2)^2+(t-4)^2=2t^2-4t+20. \]
To je parabola sa pozitivnim vodećim koeficijentom, minimum u temenu $t=-\dfrac{-4}{2\cdot2}=1.$ Tada je $|z|^2=2-4+20=18$, $|z|_{\min}=\sqrt{18}=3\sqrt{2}.$ Za $t=1$ dobijamo $z=3-3i.$

\textit{Rešenje:} $t=1,\ |z|_{\min}=3\sqrt{2},\ z=3-3i.$
\end{enumerate}
\clearpage
\section*{Slot 2. Kvadratna jednačina i Vijetove formule — rešenja}
\begin{enumerate}[label=\arabic*.]
\item Označimo $S=x_1+x_2=m$, $P=x_1x_2=3$. Realnost korena: $D=m^{2}-12\ge0$, to jest $|m|\ge2\sqrt3$.

\textbf{a)} $\dfrac{1}{x_1}+\dfrac{1}{x_2}=\dfrac{x_1+x_2}{x_1x_2}=\dfrac{m}{3}=\dfrac{5}{3}\Rightarrow m=5$. Provera: $D=25-12=13>0$.

\textbf{b)} $x_1^{3}+x_2^{3}=(x_1+x_2)^{3}-3x_1x_2(x_1+x_2)=m^{3}-9m=28$. Odatle $m^{3}-9m-28=(m-4)(m^{2}+4m+7)=0$; kvadratni činilac nema realne korene, znači $m=4$. Provera: $D=16-12=4>0$.

\textbf{c)} $\dfrac{1}{x_1^{3}}+\dfrac{1}{x_2^{3}}=\dfrac{x_1^{3}+x_2^{3}}{(x_1x_2)^{3}}=\dfrac{m^{3}-9m}{27}=6\Rightarrow m^{3}-9m=162$. Tada $m^{3}-9m-162=(m-6)(m^{2}+6m+27)=0$, drugi činilac bez realnih korena, znači $m=6$. Provera: $D=36-12=24>0$.

\emph{Rešenje:} a) $m=5$; \ b) $m=4$; \ c) $m=6$.

\item \textbf{a)} $D=(m-2)^{2}-4(m-3)=m^{2}-8m+16=(m-4)^{2}$. Jednaki koreni za $D=0$, to jest $m=4$. Tada dvostruki koren $x=\dfrac{m-2}{2}=1$.

\textbf{b)} Koren $0$ postoji kada je slobodan član jednak nuli: $m-3=0\Rightarrow m=3$. Tada jednačina $x^{2}-x=0$ sa korenima $0$ i $1$; drugi koren $1$. (Pošto je $D=(m-4)^{2}$ — potpun kvadrat, koreni su jednaki $1$ i $m-3$ za svako $m$.)

\textbf{c)} Za $x^{2}-2(m-2)x+(m^{2}-4)=0$ po Vijetovim formulama $S=2(m-2)$, $P=m^{2}-4=(m-2)(m+2)$. Diskriminanta $D=4(m-2)^{2}-4(m^{2}-4)=4(8-4m)=-16(m-2)$, znači $D\ge0\Leftrightarrow m\le2$.
(i) Različiti znaci $\Leftrightarrow P<0$ (tada je $D>0$ automatski): $(m-2)(m+2)<0\Rightarrow -2<m<2$.
(ii) Oba negativna $\Leftrightarrow D\ge0,\ P>0,\ S<0$: $P>0\Rightarrow m<-2$ ili $m>2$; $S<0\Rightarrow m<2$; zajedno sa $m\le2$ dobijamo $m<-2$.

\emph{Rešenje:} a) $m=4$, $x_1=x_2=1$; \ b) $m=3$, drugi koren $1$; \ c) različiti znaci za $-2<m<2$, oba negativna za $m<-2$.

\item \textbf{a)} Neka je $x_1=3x_2$. Po Vijetovim formulama $x_1x_2=12$, to jest $3x_2^{2}=12\Rightarrow x_2=\pm2$. Tada $x_2=2,\ x_1=6$ ili $x_2=-2,\ x_1=-6$. Zbir $m=4x_2=\pm8$. Provera: $D=64-48=16>0$.

\textbf{b)} Neka je $x_1=x_2^{2}$. Po Vijetovim formulama $x_1x_2=8$, znači $x_2^{3}=8\Rightarrow x_2=2$, tada $x_1=4$. Zbir $m=6$. Provera: $x^{2}-6x+8=(x-2)(x-4)$, $D=4>0$.

\textbf{c)} Po Vijetovim formulama $x_1+x_2=5$. Zajedno sa uslovom $2x_1-x_2=7$: sabiranjem, $3x_1=12\Rightarrow x_1=4$, tada $x_2=1$. Proizvod $m=x_1x_2=4$. Provera: $x^{2}-5x+4=(x-1)(x-4)$, $D=9>0$, i $2\cdot4-1=7$.

\emph{Rešenje:} a) $m=\pm8$; \ b) $m=6$; \ c) $m=4$.

\item \textbf{a)} Po Vijetovim formulama $x_1+x_2=m$, $x_1x_2=m^{2}-3m$. Realnost: $D=m^{2}-4(m^{2}-3m)=-3m(m-4)\ge0\Leftrightarrow 0\le m\le4$. Izraz $g(m)=x_1^{2}+x_2^{2}=(x_1+x_2)^{2}-2x_1x_2=m^{2}-2(m^{2}-3m)=-m^{2}+6m$ — parabola sa granama nadole i temenom u $m=3\in[0,4]$. Najveća vrednost $g(3)=-9+18=9$. (Za $m=3$: $D=9>0$, koreni $0$ i $3$.)

\textbf{b)} Po Vijetovim formulama $x_1+x_2=2m$, $x_1x_2=3m+4$. Realnost: $D=4m^{2}-4(3m+4)=4(m-4)(m+1)\ge0\Leftrightarrow m\le-1$ ili $m\ge4$. Izraz $g(m)=(2m)^{2}-2(3m+4)=4m^{2}-6m-8$ — parabola sa granama nagore i temenom u $m=\tfrac34$, koje leži u zabranjenom intervalu $(-1,4)$. Zato je minimum na granicama: $g(-1)=4+6-8=2$, $g(4)=64-24-8=32$. Globalni minimum $=2$ za $m=-1$ (tada $D=0$, dvostruki koren $x_1=x_2=-1$).

\emph{Rešenje:} a) $\max=9$ za $m=3$; \ b) $\min=2$ za $m=-1$.

\item \textbf{a)} Članovi progresije $a-d,\ a,\ a+d$. Zbir: $3a=15\Rightarrow a=5$. Zbir kvadrata: $3a^{2}+2d^{2}=83\Rightarrow 75+2d^{2}=83\Rightarrow d^{2}=4\Rightarrow d=\pm2$. U oba slučaja brojevi $3,5,7$.

\textbf{b)} Članovi $\dfrac{b}{q},\ b,\ bq$. Proizvod: $b^{3}=216\Rightarrow b=6$. Tada $\dfrac{6}{q}+6+6q=26\Rightarrow 6q^{2}-20q+6=0\Rightarrow 3q^{2}-10q+3=0$, $D=64$, $q=\dfrac{10\pm8}{6}$, odakle $q=3$ ili $q=\tfrac13$. Brojevi $2,6,18$.

\textbf{c)} Tačke promene znaka: $x=1$ i $x=2$. Za $x<1$: $(1-x)+(4-2x)=5-3x=5\Rightarrow x=0$ (pripada). Za $1\le x\le2$: $(x-1)+(4-2x)=3-x=5\Rightarrow x=-2$ — ne pripada. Za $x>2$: $(x-1)+(2x-4)=3x-5=5\Rightarrow x=\dfrac{10}{3}$ (pripada).

\emph{Rešenje:} a) $3,5,7$; \ b) $2,6,18$; \ c) $x\in\left\{0,\ \dfrac{10}{3}\right\}$.
\end{enumerate}
\clearpage
\section*{Slot 3. Logaritmi — rešenja}
\begin{enumerate}[label=\arabic*.]
\item
\textbf{a)} Logaritam je definisan za pozitivan argument: $x^{2}-2x>0$, to jest $x(x-2)>0$, odakle $x<0$ ili $x>2$. Znači $D_f=(-\infty,0)\cup(2,+\infty)$.

\textbf{b)} Osnova $\tfrac13<1$, pa je logaritam opadajući i pri antilogaritmovanju znak nejednačine se menja. Pošto je $-1=\log_{1/3}3$, dobijamo
\[
\log_{1/3}\left(x^{2}-2x\right) > \log_{1/3}3 \;\Rightarrow\; 0 < x^{2}-2x < 3.
\]
Levo: $x^{2}-2x>0 \Rightarrow x<0$ ili $x>2$. Desno: $x^{2}-2x-3<0 \Rightarrow (x-3)(x+1)<0 \Rightarrow -1<x<3$. Presek: $x\in(-1,0)\cup(2,3)$.

\textit{Rešenje:} a) $x\in(-\infty,0)\cup(2,+\infty)$; \quad b) $x\in(-1,0)\cup(2,3)$.

\item
\textbf{a)} Oblast definisanosti: $x>0$. Pošto je $\log(x^{3})=3\log x$, smena $t=\log x$ daje
\[
t^{2}-3t-4=0 \;\Rightarrow\; (t-4)(t+1)=0 \;\Rightarrow\; t=4 \text{ ili } t=-1,
\]
odakle $\log x=4 \Rightarrow x=10000$ i $\log x=-1 \Rightarrow x=\tfrac1{10}$.

\textbf{b)} Oblast definisanosti: $x>0,\ x\ne1$. Po formuli prelaska $\log_{x}2=\dfrac{1}{\log_{2}x}$. Smena $t=\log_{2}x\ (t\ne0)$:
\[
t+\frac{1}{t}=\frac{5}{2} \;\Rightarrow\; 2t^{2}-5t+2=0 \;\Rightarrow\; (2t-1)(t-2)=0 \;\Rightarrow\; t=2 \text{ ili } t=\tfrac12.
\]
Tada $\log_{2}x=2 \Rightarrow x=4$ i $\log_{2}x=\tfrac12 \Rightarrow x=\sqrt{2}$. Oba zadovoljavaju oblast definisanosti.

\textit{Rešenje:} a) $x=10000$ ili $x=\tfrac{1}{10}$; \quad b) $x=4$ ili $x=\sqrt{2}$.

\item Prevedimo sve na osnovu $2$ po formuli $\log_{p}q=\dfrac{\log_{2}q}{\log_{2}p}$, imajući u vidu $\log_{2}2=1$.

\textbf{a)} Pošto je $150=2\cdot3\cdot5^{2}$ i $12=2^{2}\cdot3$, to je
\[
\log_{2}150=1+a+2b,\qquad \log_{2}12=2+a,
\]
odakle $\log_{12}150=\dfrac{1+a+2b}{2+a}$.

\textbf{b)} Napišimo $0{,}24=\dfrac{6}{25}=\dfrac{2\cdot3}{5^{2}}$, tada je $\log_{2}0{,}24=1+a-2b$, a $\log_{2}5=b$, pa je
\[
\log_{5}0{,}24=\frac{\log_{2}0{,}24}{\log_{2}5}=\frac{1+a-2b}{b}.
\]

\textit{Rešenje:} a) $\dfrac{1+a+2b}{2+a}$; \quad b) $\dfrac{1+a-2b}{b}$.

\item
\textbf{a)} Oblast definisanosti: $x-1>0,\ x-1\ne1$, to jest $x>1,\ x\ne2$. Po definiciji logaritma $2x-2=(x-1)^{2}$. Pošto je $2x-2=2(x-1)$,
\[
(x-1)^{2}-2(x-1)=0 \;\Rightarrow\; (x-1)(x-3)=0 \;\Rightarrow\; x=1 \text{ ili } x=3.
\]
Koren $x=1$ ne pripada oblasti definisanosti; ostaje $x=3$ (provera: $\log_{2}4=2$).

\textbf{b)} Uslov aritmetičke progresije: $2\log_{3}\left(2^{x}-1\right)=\log_{3}2+\log_{3}\left(2^{x}+3\right)$, to jest $\left(2^{x}-1\right)^{2}=2\left(2^{x}+3\right)$. Smena $u=2^{x}>0$:
\[
u^{2}-2u+1=2u+6 \;\Rightarrow\; u^{2}-4u-5=0 \;\Rightarrow\; (u-5)(u+1)=0.
\]
Odgovara $u=5$ (koren $u=-1<0$ odbacujemo). Tada $2^{x}=5$, to jest $x=\log_{2}5$.

\textit{Rešenje:} a) $x=3$; \quad b) $x=\log_{2}5$.

\item
\textbf{a)} Oblast definisanosti: $x>0,\ y>0$. Smena $u=\log_{2}x,\ v=\log_{2}y$ daje linearni sistem $u+v=5,\ u-v=1$, odakle $u=3,\ v=2$. Znači $x=2^{3}=8,\ y=2^{2}=4$.

\textbf{b)} Oblast definisanosti: $x>0,\ y>0$. Jednačine su ekvivalentne sa $\log(xy)=\log20$ i $\log\dfrac{x}{y}=\log5$, to jest $xy=20,\ \dfrac{x}{y}=5$. Iz drugog $x=5y$, tada $5y^{2}=20 \Rightarrow y^{2}=4 \Rightarrow y=2$ (pozitivan koren), $x=10$.

\textit{Rešenje:} a) $x=8,\ y=4$; \quad b) $x=10,\ y=2$.
\end{enumerate}
\clearpage
\section*{Slot 4. Eksponencijalne jednačine i nejednačine — rešenja}
\begin{enumerate}[label=\arabic*.]
\item \textbf{a)} Neka je $t=2^{x}>0$, tada $4^{x}=t^{2}$ i $t^{2}-3t-4=0$, to jest $(t-4)(t+1)=0$. Koren $t=-1<0$ odbacujemo, ostaje $t=4$: $2^{x}=2^{2}$, $x=2$.

\textbf{b)} Neka je $t=3^{x}>0$, tada $9^{x}=t^{2}$ i $t^{2}-4t+3\le 0$, to jest $(t-1)(t-3)\le 0$, odakle $1\le t\le 3$. Znači $3^{0}\le 3^{x}\le 3^{1}$; osnova $3>1$ — funkcija raste, pa je $0\le x\le 1$.

\emph{Rešenje:} a) $x=2$; b) $x\in[0,1]$.

\item Obe jednačine su homogene drugog stepena u odnosu na $3^{x}$ i $2^{x}$. Podelimo sa $4^{x}>0$ i uvedimo $t=\left(\dfrac{3}{2}\right)^{x}>0$, pri čemu je $\dfrac{9^{x}}{4^{x}}=t^{2}$, $\dfrac{6^{x}}{4^{x}}=t$.

\textbf{a)} $4t^{2}-13t+9=0$; $D=169-144=25$, $\sqrt{D}=5$, $t=\dfrac{13\pm5}{8}$, to jest $t=\dfrac{9}{4}$ ili $t=1$. Iz $\left(\dfrac{3}{2}\right)^{x}=1$ sledi $x=0$; iz $\left(\dfrac{3}{2}\right)^{x}=\left(\dfrac{3}{2}\right)^{2}$ sledi $x=2$.

\textbf{b)} $6t^{2}-13t+6=0$; $D=25$, $t=\dfrac{13\pm5}{12}$, to jest $t=\dfrac{3}{2}$ ili $t=\dfrac{2}{3}$. Iz $\left(\dfrac{3}{2}\right)^{x}=\dfrac{3}{2}$ sledi $x=1$; iz $\left(\dfrac{3}{2}\right)^{x}=\left(\dfrac{3}{2}\right)^{-1}$ sledi $x=-1$.

\emph{Rešenje:} a) $x\in\{0,\,2\}$; b) $x\in\{-1,\,1\}$.

\item \textbf{a)} $25^{|x|}=5^{2|x|}$, pa je $x^{2}-3=2|x|$. Neka je $y=|x|\ge 0$, tada $y^{2}-2y-3=0$, $(y-3)(y+1)=0$. Koren $y=-1$ odbacujemo, ostaje $y=3$: $|x|=3$, $x=\pm3$.

\textbf{b)} $8=2^{3}$, osnova $2>1$, znači $|x^{2}-2x|\le 3$, to jest $-3\le x^{2}-2x\le 3$. Leva nejednačina $x^{2}-2x+3\ge 0$ važi za sve $x$ ($D=4-12<0$). Desna: $x^{2}-2x-3\le 0$, $(x-3)(x+1)\le 0$, $-1\le x\le 3$. Presek: $x\in[-1,3]$.

\emph{Rešenje:} a) $x\in\{-3,\,3\}$; b) $x\in[-1,3]$.

\item \textbf{a)} Oblast definisanosti: $x\ge -1$. Iz jednakosti stepena ($2>1$): $x-1=\sqrt{x+1}$, odakle dodatno $x\ge 1$. Kvadriramo: $(x-1)^{2}=x+1$, $x^{2}-3x=0$, $x(x-3)=0$. Koren $x=0$ ne zadovoljava $x\ge 1$ — strani; $x=3$ odgovara: $3-1=2=\sqrt{4}$.

\textbf{b)} Oblast definisanosti: $x\ge 0$. Neka je $t=2^{\sqrt{x}}\ge 1$, tada $4^{\sqrt{x}}=t^{2}$ i $t^{2}-5t+4=0$, $(t-1)(t-4)=0$, $t=1$ ili $t=4$. Iz $2^{\sqrt{x}}=2^{0}$: $\sqrt{x}=0$, $x=0$; iz $2^{\sqrt{x}}=2^{2}$: $\sqrt{x}=2$, $x=4$.

\emph{Rešenje:} a) $x=3$; b) $x\in\{0,\,4\}$.

\item \textbf{a)} $\log_{2}A=0\Leftrightarrow A=1$, znači $\log_{3}(9^{x}-6)=1$, odakle $9^{x}-6=3$, $9^{x}=9$, $x=1$. Provera: $9^{1}-6=3>0$ i $\log_{3}3=1>0$ — sve je definisano.

\textbf{b)} Oblast definisanosti: $25^{x}-20>0$. Po definiciji logaritma $25^{x}-20=5^{x}$. Neka je $t=5^{x}>0$, tada $t^{2}-t-20=0$, $(t-5)(t+4)=0$. Koren $t=-4$ odbacujemo, ostaje $t=5$: $5^{x}=5$, $x=1$. Provera: $25-20=5>0$, $\log_{5}5=1=x$.

\emph{Rešenje:} a) $x=1$; b) $x=1$.
\end{enumerate}
\clearpage
\section*{Slot 5. Trigonometrija na intervalu — rešenja}
\begin{enumerate}[label=\arabic*.]
\item \textbf{a)} Zamenimo $\cos 2x = 2\cos^2 x - 1$:
\[2\cos^2 x + 3\cos x + 1 = 0.\]
Neka je $t=\cos x$: $2t^2+3t+1=0$, $D=1$, $t=-\dfrac12$ ili $t=-1$.
Za $\cos x=-\dfrac12$ na $[0,2\pi)$: $x=\dfrac{2\pi}{3},\ \dfrac{4\pi}{3}$. Za $\cos x=-1$: $x=\pi$.

\textbf{b)} Pošto je $\cos 3x+\cos x = 2\cos 2x\cos x$, jednačina poprima oblik
\[\cos 2x\,(2\cos x-1)=0.\]
Iz $\cos 2x=0$: $x=\dfrac{\pi}{4}+\dfrac{k\pi}{2}$, na $(0,\pi)$ to je $\dfrac{\pi}{4},\ \dfrac{3\pi}{4}$. Iz $\cos x=\dfrac12$ na $(0,\pi)$: $x=\dfrac{\pi}{3}$.

\textbf{Rešenje.} a) $x\in\left\{\dfrac{2\pi}{3},\ \pi,\ \dfrac{4\pi}{3}\right\}$;\quad b) $x\in\left\{\dfrac{\pi}{4},\ \dfrac{\pi}{3},\ \dfrac{3\pi}{4}\right\}$.

\item \textbf{a)} Neka je $s=\sin x+\cos x$, tada $\sin x\cos x=\dfrac{s^2-1}{2}$ i $|s|\le\sqrt2$. Razlaganje zbira kubova daje
\[\sin^3 x+\cos^3 x=s\left(1-\frac{s^2-1}{2}\right)=\frac{s(3-s^2)}{2}.\]
Jednačina $\dfrac{s(3-s^2)}{2}=1$ ekvivalentna je sa $s^3-3s+2=0$, to jest $(s-1)^2(s+2)=0$. Koren $s=-2$ odbacujemo, ostaje $s=1$, to jest $\sqrt2\sin\!\left(x+\dfrac{\pi}{4}\right)=1$. Odatle $x=2k\pi$ ili $x=\dfrac{\pi}{2}+2k\pi$; na $[0,2\pi)$: $x=0$ i $x=\dfrac{\pi}{2}$.

\textbf{b)} $\sin 2x-\sin x=0\Rightarrow \sin x\,(2\cos x-1)=0$. Iz $\sin x=0$: $x=0,\ \pi$. Iz $\cos x=\dfrac12$: $x=\dfrac{\pi}{3},\ \dfrac{5\pi}{3}$.

\textbf{Rešenje.} a) $x\in\left\{0,\ \dfrac{\pi}{2}\right\}$;\quad b) $x\in\left\{0,\ \dfrac{\pi}{3},\ \pi,\ \dfrac{5\pi}{3}\right\}$.

\item Oblast definisanosti u oba zadatka: $\sin x\ne 0$ i $\cos x\ne 0$, to jest na $[0,2\pi)$ isključuju se tačke $0,\ \dfrac{\pi}{2},\ \pi,\ \dfrac{3\pi}{2}$.

\textbf{a)} $\operatorname{tg} x+\operatorname{ctg} x=\dfrac{1}{\sin x\cos x}=\dfrac{2}{\sin 2x}$. Jednačina $\dfrac{2}{\sin 2x}=2$ daje $\sin 2x=1$, $x=\dfrac{\pi}{4}+k\pi$; na $[0,2\pi)$: $x=\dfrac{\pi}{4},\ \dfrac{5\pi}{4}$ (obe u oblasti definisanosti).

\textbf{b)} Zamenom $\operatorname{ctg} x=\dfrac{1}{\operatorname{tg} x}$, dobijamo $\operatorname{tg}^2 x=3$, to jest $\operatorname{tg} x=\pm\sqrt3$. Iz $\operatorname{tg} x=\sqrt3$: $x=\dfrac{\pi}{3},\ \dfrac{4\pi}{3}$; iz $\operatorname{tg} x=-\sqrt3$: $x=\dfrac{2\pi}{3},\ \dfrac{5\pi}{3}$. Sve tačke u oblasti definisanosti.

\textbf{Rešenje.} a) $x\in\left\{\dfrac{\pi}{4},\ \dfrac{5\pi}{4}\right\}$;\quad b) $x\in\left\{\dfrac{\pi}{3},\ \dfrac{2\pi}{3},\ \dfrac{4\pi}{3},\ \dfrac{5\pi}{3}\right\}$.

\item Iz $\sin\!\left(x-\dfrac{\pi}{6}\right)=\dfrac{\sqrt3}{2}$ dobijamo dve serije:
\[x-\frac{\pi}{6}=\frac{\pi}{3}+2k\pi\quad\text{ili}\quad x-\frac{\pi}{6}=\frac{2\pi}{3}+2k\pi,\]
to jest $x=\dfrac{\pi}{2}+2k\pi$ ili $x=\dfrac{5\pi}{6}+2k\pi$. Na $[0,2\pi)$ ostaju $x=\dfrac{\pi}{2}$ i $x=\dfrac{5\pi}{6}$.

Odabir korena koji pripadaju $\left(\dfrac{2\pi}{3},2\pi\right)$: $\dfrac{\pi}{2}\approx1{,}57<\dfrac{2\pi}{3}\approx2{,}09$ ne pripada, a $\dfrac{5\pi}{6}\approx2{,}62$ pripada.

\textbf{Rešenje.} Na $[0,2\pi)$: $x\in\left\{\dfrac{\pi}{2},\ \dfrac{5\pi}{6}\right\}$; u $\left(\dfrac{2\pi}{3},2\pi\right)$: samo $x=\dfrac{5\pi}{6}$.

\item Koristeći $\cos 2x=2\cos^2 x-1$, rastavimo funkciju na činioce:
\[f(x)=2\cos x-(2\cos^2 x-1)-1=2\cos x\,(1-\cos x).\]

\textbf{a)} $f(x)=0\iff\cos x=0$ ili $\cos x=1$. Na $(-\pi,\pi]$ iz $\cos x=0$ dobijamo $x=\pm\dfrac{\pi}{2}$, iz $\cos x=1$ dobijamo $x=0$. Nule: $-\dfrac{\pi}{2},\ 0,\ \dfrac{\pi}{2}$.

Činilac $1-\cos x\ge 0$ svuda i jednak je nuli samo za $x=0$, pa za $x\ne 0$ znak $f$ se poklapa sa znakom $\cos x$.

\textbf{b)} $f(x)>0\iff\cos x>0$ i $x\ne 0$. Na $(-\pi,\pi]$ to je $x\in\left(-\dfrac{\pi}{2},0\right)\cup\left(0,\dfrac{\pi}{2}\right)$.

\textbf{c)} $f(x)<0\iff\cos x<0$, to jest $x\in\left(-\pi,-\dfrac{\pi}{2}\right)\cup\left(\dfrac{\pi}{2},\pi\right]$ (u tački $x=\pi$ imamo $f(\pi)=-4<0$).

\textbf{Rešenje.} a) $x\in\left\{-\dfrac{\pi}{2},\ 0,\ \dfrac{\pi}{2}\right\}$;\quad b) $x\in\left(-\dfrac{\pi}{2},0\right)\cup\left(0,\dfrac{\pi}{2}\right)$;\quad c) $x\in\left(-\pi,-\dfrac{\pi}{2}\right)\cup\left(\dfrac{\pi}{2},\pi\right]$.
\end{enumerate}
\clearpage
\section*{Slot 6. Vektori i analitička geometrija — rešenja}
\begin{enumerate}[label=\arabic*.]

\item Po uslovu $m\cdot n=|m|\,|n|\cos\dfrac{\pi}{3}=2\cdot 3\cdot\dfrac12=3$, $|m|^2=4$, $|n|^2=9$.

a) $|a|^2=|2m-n|^2=4|m|^2-4(m\cdot n)+|n|^2=16-12+9=13$, znači $|a|=\sqrt{13}$.

b) Površina je jednaka $|a\times b|$. Pošto je $m\times m=n\times n=0$ i $n\times m=-(m\times n)$,
\[a\times b=(2m-n)\times(m+3n)=6\,(m\times n)-(n\times m)=7\,(m\times n).\]
Pri tome je $|m\times n|=|m|\,|n|\sin\dfrac{\pi}{3}=2\cdot3\cdot\dfrac{\sqrt3}{2}=3\sqrt3$, pa je $S=7\cdot 3\sqrt3=21\sqrt3$.

v) Uslov $(m+tn)\cdot(2m-n)=0$ daje
\[2|m|^2-(m\cdot n)+2t(m\cdot n)-t|n|^2=8-3+6t-9t=5-3t=0,\]
odakle $t=\dfrac{5}{3}$.

g) Iz $m=xa+yb$ dobijamo $(2x+y)m+(-x+3y)n=1\cdot m+0\cdot n$. Pošto su $m,n$ nezavisni, $\begin{cases}2x+y=1\\ -x+3y=0\end{cases}$. Iz drugog $x=3y$, tada $7y=1$, $y=\dfrac17$, $x=\dfrac37$.

\textit{Rešenje:} a) $|a|=\sqrt{13}$; b) $S=21\sqrt3$; v) $t=\dfrac{5}{3}$; g) $m=\dfrac37a+\dfrac17b$.

\item Vektori: $\overrightarrow{AB}=(-2,1,1)$, $\overrightarrow{AC}=(-1,-1,2)$.

a) $\overrightarrow{AB}\cdot\overrightarrow{AC}=2-1+2=3$, $|\overrightarrow{AB}|=|\overrightarrow{AC}|=\sqrt6$. Tada
\[\cos\angle BAC=\dfrac{3}{\sqrt6\cdot\sqrt6}=\dfrac12,\quad\angle BAC=60^\circ.\]

b) $\overrightarrow{AB}\times\overrightarrow{AC}=(3,3,3)$, njegova dužina $\sqrt{27}=3\sqrt3$. Znači $S=\dfrac12\cdot 3\sqrt3=\dfrac{3\sqrt3}{2}$.

v) $T=\dfrac{A+B+C}{3}=\left(\dfrac{1-1+0}{3},\dfrac{1+2+0}{3},\dfrac{1+2+3}{3}\right)=(0,1,2)$.

\textit{Rešenje:} a) $\angle BAC=60^\circ$; b) $S=\dfrac{3\sqrt3}{2}$; v) $T(0,1,2)$.

\item a) Ako su dijagonale paralelograma jednake $d_1$ i $d_2$, tada je njegova površina $S=\dfrac12\,|d_1\times d_2|$. Imamo $d_1\times d_2=(6,6,-6)$, $|d_1\times d_2|=\sqrt{108}=6\sqrt3$, pa je $S=\dfrac12\cdot 6\sqrt3=3\sqrt3$.

b) Vektori su komplanarni ako i samo ako je njihov mešoviti proizvod jednak nuli:
\[\begin{vmatrix}1&2&-1\\ 2&-1&3\\ \lambda&1&2\end{vmatrix}=1(-2-3)-2(4-3\lambda)+(-1)(2+\lambda)=5\lambda-15.\]
Iz $5\lambda-15=0$ dobijamo $\lambda=3$.

v) Kolinearnost znači proporcionalnost koordinata: $\dfrac{2}{-4}=\dfrac{-3}{6}=\dfrac{k}{10}$. Zajednička vrednost prva dva odnosa jednaka je $-\dfrac12$, pa je $\dfrac{k}{10}=-\dfrac12$, odakle $k=-5$.

\textit{Rešenje:} a) $S=3\sqrt3$; b) $\lambda=3$; v) $k=-5$.

\item a) $p\cdot q=|p|\,|q|\cos\dfrac{\pi}{4}=1\cdot\sqrt2\cdot\dfrac{\sqrt2}{2}=1$. Tada $|p+q|^2=|p|^2+2(p\cdot q)+|q|^2=1+2+2=5$, znači $|p+q|=\sqrt5$. Uslov $(p+tq)\cdot q=0$ daje $p\cdot q+t|q|^2=1+2t=0$, odakle $t=-\dfrac12$.

b) $a\cdot b=3\alpha-\alpha-4=2\alpha-4=0$, pa je $\alpha=2$.

v) $|u|^2=(\alpha-1)^2+9$, $|v|^2=(\alpha+1)^2+5$. Iz $|u|^2=|v|^2$:
\[(\alpha-1)^2+9=(\alpha+1)^2+5\ \Rightarrow\ 4-4\alpha=0,\]
odakle $\alpha=1$.

\textit{Rešenje:} a) $|p+q|=\sqrt5$, $t=-\dfrac12$; b) $\alpha=2$; v) $\alpha=1$.

\item a) Visina iz $A$ je normalna na $BC$: $\overrightarrow{BC}=(-4,4)$, jednačina $-4(x-1)+4(y-1)=0$, to jest $y=x$. Visina iz $B$ je normalna na $AC$: $\overrightarrow{AC}=(2,4)$, jednačina $2(x-7)+4(y-1)=0$, to jest $x+2y=9$. Rešavajući $\begin{cases}y=x\\ x+2y=9\end{cases}$, dobijamo $x=3$, $y=3$. Ortocentar $H(3,3)$.

b) Centar $O(x,y)$ je jednako udaljen od temena: iz $|OA|^2=|OB|^2$ sledi $(x-1)^2=(x-7)^2$, to jest $x=4$; iz $|OA|^2=|OC|^2$: $9+(y-1)^2=1+(y-5)^2$, to jest $8y=16$, $y=2$. Tada $r^2=(4-1)^2+(2-1)^2=10$, $r=\sqrt{10}$. Jednačina: $(x-4)^2+(y-2)^2=10$.

v) Stranice: $\overrightarrow{PQ}=(3,4)$, $|PQ|=5$; $\overrightarrow{QR}=(5,0)$, $|QR|=5$; $\overrightarrow{RS}=(-3,-4)$, $|RS|=5$; $\overrightarrow{SP}=(-5,0)$, $|SP|=5$. Sve četiri stranice su jednake $5$. Osim toga $\overrightarrow{PQ}=-\overrightarrow{RS}$, znači $PQRS$ je paralelogram; paralelogram sa jednakim stranicama je romb. Dodatno: dijagonale $\overrightarrow{PR}=(8,4)$ i $\overrightarrow{QS}=(2,-4)$ su ortogonalne, jer je $8\cdot2+4\cdot(-4)=0$.

\textit{Rešenje:} a) $H(3,3)$; b) $O(4,2)$, $r=\sqrt{10}$, $(x-4)^2+(y-2)^2=10$; v) $PQRS$ je romb.

\end{enumerate}
\clearpage
\section*{Slot 7. Planimetrija — rešenja}
\begin{enumerate}[label=\arabic*.]

\item Spustimo iz krajeva manje osnovice dve visine na veću osnovicu; one odsecaju od $a$ dva horizontalna odsečka $p$ i $q$, pri čemu je $a-c=p+q=12$. Iz pravouglih trouglova $p=\dfrac{h}{\tan 30^\circ}=h\sqrt3$, $q=\dfrac{h}{\tan 45^\circ}=h$, pa je
\[ h\sqrt3+h=12 \Rightarrow h=\frac{12}{\sqrt3+1}=6(\sqrt3-1). \]
a) Krakovi preko sinusa istih uglova:
\[ b_{30}=\frac{h}{\sin 30^\circ}=2h=12(\sqrt3-1),\qquad b_{45}=\frac{h}{\sin 45^\circ}=h\sqrt2=6(\sqrt6-\sqrt2). \]
b) Površina $P=\dfrac{a+c}{2}\,h=\dfrac{20+8}{2}\cdot 6(\sqrt3-1)=84(\sqrt3-1).$

\textbf{Rešenje:} a) $h=6(\sqrt3-1)$, krakovi $12(\sqrt3-1)$ i $6(\sqrt6-\sqrt2)$; \; b) $P=84(\sqrt3-1)$.

\item a) Uslov upisane kružnice: $a+c=b+d$, a za jednakokraki trapez $b=d$, pa je $b=\dfrac{a+c}{2}=\dfrac{18+8}{2}=13$. Projekcija kraka na veću osnovicu jednaka je $\dfrac{a-c}{2}=5$, tada $h=\sqrt{b^2-5^2}=\sqrt{169-25}=12$. Kružnica dodiruje obe osnovice, znači njen prečnik je jednak visini: $r=\dfrac{h}{2}=6$. Površina $P=\dfrac{a+c}{2}\,h=13\cdot 12=156$.

b) Postavimo u koordinate $A(-7,0),B(7,0),C(3,h),D(-3,h)$. Uslov $\overrightarrow{AC}\perp\overrightarrow{BD}$ daje $\overrightarrow{AC}\cdot\overrightarrow{BD}=(10)(-10)+h^2=0$, to jest $h^2=100$, $h=10=\dfrac{a+c}{2}$. Površina $P=\dfrac{a+c}{2}\,h=10\cdot 10=100$.

\textbf{Rešenje:} a) $b=13$, $r=6$, $P=156$; \; b) $h=10$, $P=100$.

\item Dijagonale romba su normalne i polove se, pa je $d_1^2+d_2^2=4a^2$.

a) $d_1^2+d_2^2=4\cdot 13^2=676$. Iz $d_1-d_2=14$: $d_1d_2=\dfrac{676-196}{2}=240$, tada $(d_1+d_2)^2=676+480=1156$, $d_1+d_2=34$. Odatle $d_1=24,\ d_2=10$. Površina $P=\dfrac{d_1d_2}{2}=120$; visina $h=\dfrac{P}{a}=\dfrac{120}{13}$, poluprečnik $r=\dfrac{h}{2}=\dfrac{60}{13}$.

b) Neka je $d_1=3k,\ d_2=4k$. Iz $\dfrac{3k\cdot 4k}{2}=6k^2=96$ dobijamo $k=4$, znači $d_1=12,\ d_2=16$. Stranica $a=\sqrt{6^2+8^2}=10$; $h=\dfrac{P}{a}=\dfrac{96}{10}=9{,}6$, poluprečnik $r=\dfrac{h}{2}=\dfrac{24}{5}=4{,}8$.

\textbf{Rešenje:} a) $d_1=24,\ d_2=10$, $P=120$, $r=\dfrac{60}{13}$; \; b) $d_1=12,\ d_2=16$, $a=10$, $r=\dfrac{24}{5}$.

\item a) Neka su stranice $x-d,\ x,\ x+d$. Zbir $3x=30$, znači $x=10$. Ugao $120^\circ$ leži naspram najveće stranice $x+d$. Po kosinusnoj teoremi sa $\cos 120^\circ=-\tfrac12$:
\[ (10+d)^2=(10-d)^2+10^2-2(10-d)\cdot 10\cdot\left(-\tfrac12\right), \]
odakle $40d=200-10d$, $d=4$. Stranice $6,\ 10,\ 14$ (provera: $14^2=6^2+10^2+6\cdot 10=196$). Površina
\[ S=\frac12\cdot 6\cdot 10\cdot\sin 120^\circ=30\cdot\frac{\sqrt3}{2}=15\sqrt3. \]
b) Poluobim $s=15$. Tada $r=\dfrac{S}{s}=\dfrac{15\sqrt3}{15}=\sqrt3$, a po sinusnoj teoremi
\[ R=\frac{14}{2\sin 120^\circ}=\frac{14}{\sqrt3}=\frac{14\sqrt3}{3}. \]

\textbf{Rešenje:} a) stranice $6,\ 10,\ 14$, $S=15\sqrt3$; \; b) $r=\sqrt3$, $R=\dfrac{14\sqrt3}{3}$.

\item a) Pri spoljašnjem dodiru rastojanje između centara $d=R_1+R_2=13$. Poluprečnici u tačkama dodira su normalni na tangentu, pa je odsečak tangente — kateta pravouglog trapeza sa osnovicama $R_1,R_2$ i krakom $d$:
\[ \ell=\sqrt{d^2-(R_1-R_2)^2}=\sqrt{13^2-5^2}=\sqrt{144}=12. \]
Površina trapeza (paralelne stranice $9$ i $4$, visina $\ell=12$): $P=\dfrac{9+4}{2}\cdot 12=78$.

b) Za tri kružnice koje dodiruju jednu pravu i jedna drugu po parovima, projekcije tačaka dodira se sabiraju: $2\sqrt{R_1\rho}+2\sqrt{R_2\rho}=2\sqrt{R_1R_2}$, to jest
\[ \frac{1}{\sqrt{\rho}}=\frac{1}{\sqrt{R_1}}+\frac{1}{\sqrt{R_2}}=\frac13+\frac12=\frac56\;\Rightarrow\;\sqrt{\rho}=\frac65\;\Rightarrow\;\rho=\frac{36}{25}. \]

\textbf{Rešenje:} a) dužina tangente $=12$, površina trapeza $=78$; \; b) $\rho=\dfrac{36}{25}$.

\item a) Normala iz centra polovi tetivu: polovina tetive $=\sqrt{R^2-3^2}=\sqrt{36-9}=3\sqrt3$, pa je tetiva jednaka $6\sqrt3$. Polovina centralnog ugla $\alpha$ zadovoljava $\cos\dfrac{\alpha}{2}=\dfrac{3}{6}=\dfrac12$, odakle $\dfrac{\alpha}{2}=60^\circ$, $\alpha=120^\circ$. Periferijski ugao je dvostruko manji od centralnog: $\beta=60^\circ$.

b) Manji segment odgovara centralnom uglu $\theta=120^\circ=\dfrac{2\pi}{3}$. Površina segmenta:
\[ P_{\text{seg}}=\frac{R^2}{2}\bigl(\theta-\sin\theta\bigr)=\frac{36}{2}\left(\frac{2\pi}{3}-\frac{\sqrt3}{2}\right)=12\pi-9\sqrt3. \]
Udeo u površini kruga $36\pi$:
\[ \frac{12\pi-9\sqrt3}{36\pi}=\frac13-\frac{\sqrt3}{4\pi}\approx 0{,}1955=19{,}55\,\%. \]

\textbf{Rešenje:} a) tetiva $=6\sqrt3$, centralni ugao $120^\circ$, periferijski ugao $60^\circ$; \; b) $\dfrac{1}{3}-\dfrac{\sqrt3}{4\pi}\approx 19{,}55\,\%$.

\end{enumerate}
\clearpage
\section*{Slot 8. Stereometrija — rešenja}
\begin{enumerate}[label=\arabic*.]
\item \textbf{a)} Uslov progresije: $2H=r+s$, a iz osnog preseka $s^2=r^2+H^2$. Za $r=3$ imamo $s=2H-3$, pa je
\[(2H-3)^2=9+H^2\Rightarrow 3H^2-12H=0\Rightarrow H=4,\quad s=5.\]
Progresija $3,4,5$ sa razlikom $1$. Tada $V=\frac{1}{3}\pi r^2 H=\frac{1}{3}\pi\cdot 9\cdot 4=12\pi$ i $P=\pi r^2+\pi r s=9\pi+15\pi=24\pi$.
\textbf{b)} Osni presek — jednakokraki trougao sa osnovicom $2r=6$ i kracima $5$ (visina $4$). Poluprečnik upisane kružnice $\rho=\dfrac{S}{p}=\dfrac{12}{8}=\dfrac{3}{2}$.
\textbf{v)} $V_{\text{lopta}}=\dfrac{4}{3}\pi\rho^3=\dfrac{4}{3}\pi\cdot\dfrac{27}{8}=\dfrac{9}{2}\pi$, odakle $\dfrac{V_{\text{lopta}}}{V_{\text{kupa}}}=\dfrac{9\pi/2}{12\pi}=\dfrac{3}{8}$.
\emph{Rešenje:} a) $H=4,\ s=5,\ V=12\pi,\ P=24\pi$; b) $\rho=\dfrac32$; v) $\dfrac38$.

\item \textbf{a)} Horizontalni presek piramide na visini $z$ — kvadrat sa stranicom $a\cdot\dfrac{H-z}{H}$. Gornja strana kocke na visini $z=b$ ima stranicu $b$, jednaku stranici preseka:
\[b=a\cdot\frac{H-b}{H}\Rightarrow b(a+H)=aH\Rightarrow b=\frac{aH}{a+H}.\]
\textbf{b)} Za $a=6,\ H=12$: $b=\dfrac{72}{18}=4$, $V_{\text{kocka}}=4^3=64$, $P_{\text{kocka}}=6\cdot 16=96$.
\textbf{v)} $V_{\text{pir}}=\dfrac{1}{3}a^2H=\dfrac{1}{3}\cdot 36\cdot 12=144$, pa je $\dfrac{V_{\text{kocka}}}{V_{\text{pir}}}=\dfrac{64}{144}=\dfrac{4}{9}$.
\emph{Rešenje:} a) $b=\dfrac{aH}{a+H}$; b) $b=4,\ V=64,\ P=96$; v) $\dfrac49$.

\item Uzmimo $\ell$ za osu $Oy$. Trougao: prav ugao $A=(2,0)$, $B=(2,4)$, $C=(5,0)$.
\textbf{a)} Na visini $y\in[0,4]$ figura zauzima $x$ od $2$ do $R(y)=2+3\bigl(1-\tfrac{y}{4}\bigr)$; pri rotaciji — prsten sa spoljašnjim poluprečnikom $R(y)$ i unutrašnjim $2$. Dobija se valjak poluprečnika $5$ visine $4$ sa koničnim udubljenjem. Po metodi preseka
\[V=\pi\int_0^4\Bigl[9\bigl(1-\tfrac{y}{4}\bigr)^2+12\bigl(1-\tfrac{y}{4}\bigr)\Bigr]dy=\pi(12+24)=36\pi.\]
Provera po Guldinu: $\bar x=\tfrac{2+2+5}{3}=3$, $S=6$, $V=2\pi\bar x S=2\pi\cdot3\cdot6=36\pi$.
\textbf{b)} Po Guldinovoj teoremi za površinu $P=2\pi\sum\bar x_i L_i$: kateta $AB$ ($L=4,\ \bar x=2$) daje $16\pi$; kateta $AC$ ($L=3,\ \bar x=\tfrac72$) daje $21\pi$; hipotenuza $BC$ ($L=5,\ \bar x=\tfrac72$) daje $35\pi$. Ukupno $P=72\pi$.
\emph{Rešenje:} a) $V=36\pi$; b) $P=72\pi$.

\item \textbf{a)} U osnom preseku visina valjka opada linearno od $H=9$ za $x=0$ do $0$ za $x=R=6$: $h=9\bigl(1-\tfrac{x}{6}\bigr)$. Tada
\[V(x)=\pi x^2 h=\pi\Bigl(9x^2-\tfrac{3}{2}x^3\Bigr),\qquad x\in(0;6).\]
\textbf{b)} $V'(x)=\pi\bigl(18x-\tfrac{9}{2}x^2\bigr)=\dfrac{9\pi}{2}x(4-x)$: na $(0;6)$ maksimum za $x=4$. Tada $h=3$, $V_{\max}=\pi\cdot16\cdot3=48\pi$.
\textbf{v)} $V_{\text{kupa}}=\dfrac{1}{3}\pi R^2H=108\pi$, pa je $\dfrac{V_{\max}}{V_{\text{kupa}}}=\dfrac{48}{108}=\dfrac{4}{9}$.
\emph{Rešenje:} a) $h=9\bigl(1-\tfrac{x}{6}\bigr),\ V(x)=\pi\bigl(9x^2-\tfrac32x^3\bigr),\ x\in(0;6)$; b) $x=4,\ V_{\max}=48\pi$; v) $\dfrac49$.

\item \textbf{a)} Osni presek — jednakokraki trapez sa osnovicama $2R,2r$ i upisanom kružnicom poluprečnika $\rho$, pri čemu je $h=2\rho$. Za tangentni trapez $2R+2r=2s$, to jest $s=R+r=5$. Iz $s^2=h^2+(R-r)^2$ dobijamo $h^2=(R+r)^2-(R-r)^2=4Rr$, znači $h=2\sqrt{Rr}=4$ i $\rho=2$.
\textbf{b)} $V_{\text{lopta}}=\dfrac{4}{3}\pi\rho^3=\dfrac{32}{3}\pi$, $V_{\text{kupa}}=\dfrac{\pi h}{3}(R^2+Rr+r^2)=\dfrac{4\pi}{3}\cdot21=28\pi$; odnos $\dfrac{32/3}{28}=\dfrac{8}{21}$.
\textbf{v)} $P_{\text{lopta}}=4\pi\rho^2=16\pi$, $P_{\text{kupa}}=\pi(R+r)s+\pi R^2+\pi r^2=25\pi+17\pi=42\pi$; odnos $\dfrac{16}{42}=\dfrac{8}{21}$. Poklapa se sa odnosom zapremina, jer za telo sa upisanom loptom važi $V=\tfrac{\rho}{3}P$.
\emph{Rešenje:} a) $\rho=2,\ h=4,\ s=5$; b) $V_{\text{lopta}}=\tfrac{32}{3}\pi,\ V_{\text{kupa}}=28\pi,\ \tfrac{8}{21}$; v) $P_{\text{lopta}}=16\pi,\ P_{\text{kupa}}=42\pi,\ \tfrac{8}{21}$.

\item \textbf{a)} Presek kroz $S$ i sredine $M,N$ naspramnih stranica osnove — jednakokraki trougao $SMN$ sa $MN=a$ i visinom $SO=H$; $SM,SN$ — apoteme strana, a ugao između strana jednak je $\angle MSN=90^\circ$. Znači $SMN$ je pravougli jednakokraki, visina iz pravog ugla jednaka je polovini hipotenuze: $H=\dfrac{a}{2}=3$. Zapremina $V=\dfrac{1}{3}a^2H=\dfrac{1}{3}\cdot36\cdot3=36$.
\textbf{b)} Apotema $m=\sqrt{H^2+\bigl(\tfrac{a}{2}\bigr)^2}=\sqrt{9+9}=3\sqrt2$. Bočna površina $=4\cdot\tfrac12 a m=2am=36\sqrt2$, osnova $36$, ukupna $P=36+36\sqrt2$.
\textbf{v)} $\tan\varphi=\dfrac{H}{a/2}=\dfrac{3}{3}=1\Rightarrow\varphi=45^\circ$.
\emph{Rešenje:} a) $H=3,\ V=36$; b) $m=3\sqrt2,\ P=36+36\sqrt2$; v) $45^\circ$.

\item \textbf{a)} $P=V$: $6a^2=a^3\Rightarrow a=6$.
\textbf{b)} Tetraedar kod temena $A$ ima tri uzajamno normalne ivice dužine $a$, $V_{\text{tet}}=\dfrac{1}{6}a^3$. Naspramna strana — jednakostraničan trougao sa stranicom $a\sqrt2$, površina $S=\dfrac{\sqrt3}{4}(a\sqrt2)^2=\dfrac{\sqrt3}{2}a^2$. Tada
\[d=\frac{3V_{\text{tet}}}{S}=\frac{\tfrac12 a^3}{\tfrac{\sqrt3}{2}a^2}=\frac{a}{\sqrt3}=\frac{6}{\sqrt3}=2\sqrt3.\]
\textbf{v)} $r=\dfrac{a}{2}$, $R=\dfrac{a\sqrt3}{2}$, $\dfrac{r}{R}=\dfrac{1}{\sqrt3}$, pa je $\dfrac{V_{\text{up}}}{V_{\text{op}}}=\Bigl(\dfrac{1}{\sqrt3}\Bigr)^3=\dfrac{1}{3\sqrt3}=\dfrac{\sqrt3}{9}$.
\emph{Rešenje:} a) $a=6$; b) $d=2\sqrt3$; v) $\dfrac{\sqrt3}{9}$.
\end{enumerate}
\clearpage
\section*{Slot 9. Ispitivanje funkcije i integral — rešenja}
\begin{enumerate}[label=\arabic*.]
  \item \textbf{a)} Imenilac je jednak nuli za $x=2$, pa je $D=\mathbb{R}\setminus\{2\}$. Deljenjem: $f(x)=\dfrac{x^2-3}{x-2}=x+2+\dfrac{1}{x-2}$. Pošto je $\lim_{x\to 2}|f(x)|=\infty$, prava $x=2$ je vertikalna asimptota; pošto $\dfrac{1}{x-2}\to 0$ za $x\to\pm\infty$, kosa asimptota $y=x+2$.

  \textbf{b)} $f'(x)=1-\dfrac{1}{(x-2)^2}=\dfrac{(x-1)(x-3)}{(x-2)^2}$. Znak $f'$ se poklapa sa znakom $(x-1)(x-3)$: $f'>0$ na $(-\infty,1)$ i $(3,+\infty)$ (rastenje), $f'<0$ na $(1,2)$ i $(2,3)$ (opadanje). U $x=1$ — maksimum $f(1)=2$, u $x=3$ — minimum $f(3)=6$.

  \textbf{c)} Po razlaganju $\displaystyle\int_{3}^{4}\Bigl(x+2+\frac{1}{x-2}\Bigr)dx=\Bigl[\frac{x^2}{2}+2x+\ln(x-2)\Bigr]_{3}^{4}=(16+\ln 2)-\frac{21}{2}=\frac{11}{2}+\ln 2$.

  \textbf{Rešenje:} a) $D=\mathbb{R}\setminus\{2\}$, vert. $x=2$, kosa $y=x+2$; b) maks $f(1)=2$, min $f(3)=6$; c) $\dfrac{11}{2}+\ln 2$.

  \item \textbf{a)} $f'(x)=2x e^{-x}-x^2 e^{-x}=x(2-x)e^{-x}$. Pošto je $e^{-x}>0$, znak $f'$ jednak je znaku $x(2-x)$: $f'>0$ na $(0,2)$, $f'<0$ na $(-\infty,0)$ i $(2,+\infty)$. Znači $x=0$ — minimum $f(0)=0$, $x=2$ — maksimum $f(2)=\dfrac{4}{e^2}$.

  \textbf{b)} Eksponencijalna funkcija raste brže od bilo kog stepena: $\lim_{x\to+\infty}\dfrac{x^2}{e^x}=0$. Dalje $\dfrac{f(x)}{x}=x e^{-x}\to 0$, $\dfrac{f(x)}{x^2}=e^{-x}\to 0$. Pošto $f(x)\to 0$, horizontalna asimptota $y=0$.

  \textbf{c)} Dvaput parcijalnom integracijom dobijamo primitivnu funkciju $F(x)=-(x^2+2x+2)e^{-x}$. Tada $\displaystyle\int_{0}^{1}=F(1)-F(0)=-5e^{-1}+2=2-\dfrac{5}{e}$.

  \textbf{Rešenje:} a) min $f(0)=0$, maks $f(2)=\dfrac{4}{e^2}$; b) sve tri granične vrednosti $=0$, asimptota $y=0$; c) $2-\dfrac{5}{e}$.

  \item \textbf{a)} $g'(x)=3(x^2-4)^2\cdot 2x=6x(x-2)^2(x+2)^2$. Nule: $x=0$ (višestrukost $1$) i $x=\pm 2$ (višestrukost $2$). Činioci $(x\mp 2)^2$ su nenegativni, pa je znak $g'$ određen činiocem $x$: $g'<0$ za $x<0$ i $g'>0$ za $x>0$. U $x=0$ funkcija prelazi iz opadanja u rastenje — minimum $g(0)=-64$; u tačkama $x=\pm 2$ znak $g'$ se ne menja, pa nema ekstremne vrednosti (horizontalne tangente-prevoji).

  \textbf{b)} Smena $u=x^2-4$, $x\,dx=\tfrac12 du$; granice $x=2\Rightarrow u=0$, $x=3\Rightarrow u=5$:
  \[
  \int_{2}^{3} x(x^2-4)^2dx=\frac12\int_{0}^{5} u^2\,du=\frac12\cdot\frac{u^3}{3}\Big|_{0}^{5}=\frac{125}{6}.
  \]

  \textbf{c)} $h(x)=x(x-2)(x+2)$ ima nule $-2,0,2$; na $(-2,0)$ $h>0$, na $(0,2)$ $h<0$. Računamo po modulu:
  \[
  \int_{-2}^{0}(x^3-4x)dx=4,\qquad \int_{0}^{2}(x^3-4x)dx=-4,\qquad S=|4|+|-4|=8.
  \]

  \textbf{Rešenje:} a) $g'=0$ za $x=-2,0,2$, ekstremna vrednost samo u $x=0$ (min $-64$); b) $\dfrac{125}{6}$; c) $S=8$.

  \item \textbf{a)} $D=\mathbb{R}\setminus\{-1\}$. Deljenjem $f(x)=\dfrac{x^2+2}{x+1}=x-1+\dfrac{3}{x+1}$: vertikalna asimptota $x=-1$, kosa $y=x-1$.

  \textbf{b)} $f'(x)=1-\dfrac{3}{(x+1)^2}$. U tački $x_0=0$: $f(0)=2$, $f'(0)=-2$. Tangenta $y=-2x+2$; koeficijent pravca normale $-\dfrac{1}{f'(0)}=\dfrac12$, normala $y=\dfrac12 x+2$.

  \textbf{c)} $\dfrac{f(x)}{x}=\dfrac{x^2+2}{x^2+x}\to 1$ za $x\to+\infty$ (odnos vodećih koeficijenata).

  \textbf{Rešenje:} a) $D=\mathbb{R}\setminus\{-1\}$, vert. $x=-1$, kosa $y=x-1$; b) tangenta $y=-2x+2$, normala $y=\dfrac12 x+2$; c) $1$.

  \item \textbf{a)} $D=\mathbb{R}\setminus\{1\}$. $f(x)=\dfrac{3x-2}{x-1}=3+\dfrac{1}{x-1}$: za $x\to 1$ moduo teži ka $\infty$ — vertikalna asimptota $x=1$; za $x\to\pm\infty$ sabirak $\to 0$ — horizontalna asimptota $y=3$.

  \textbf{b)} Sa $t=2x$: $\dfrac{e^{2x}-1}{x}=2\cdot\dfrac{e^{2x}-1}{2x}\to 2$. Sa $t=\dfrac{3}{x}\to 0$: $x\sin\dfrac{3}{x}=3\cdot\dfrac{\sin t}{t}\to 3$.

  \textbf{c)} Množeći konjugovanim: $\sqrt{x^2+x}-x=\dfrac{x}{\sqrt{x^2+x}+x}=\dfrac{1}{\sqrt{1+1/x}+1}\to\dfrac12$.

  \textbf{Rešenje:} a) $D=\mathbb{R}\setminus\{1\}$, vert. $x=1$, horiz. $y=3$; b) $2$ i $3$; c) $\dfrac12$.

  \item \textbf{a)} Parcijalnom integracijom $u=\ln x$, $dv=x\,dx$:
  \[
  \int_{1}^{e} x\ln x\,dx=\frac{x^2}{2}\ln x\Big|_{1}^{e}-\int_{1}^{e}\frac{x}{2}\,dx=\frac{e^2}{2}-\Bigl(\frac{e^2}{4}-\frac14\Bigr)=\frac{e^2+1}{4}.
  \]

  \textbf{b)} Smena $u=1+x^2$, $x\,dx=\tfrac12 du$; granice $1\to 2$:
  \[
  \int_{0}^{1} x\sqrt{1+x^2}\,dx=\frac12\int_{1}^{2} u^{1/2}\,du=\frac13 u^{3/2}\Big|_{1}^{2}=\frac{2\sqrt2-1}{3}.
  \]

  \textbf{c)} Smena $t=\ln x$: $\displaystyle\int_{1}^{e}\frac{\ln x}{x}\,dx=\int_{0}^{1} t\,dt=\frac12$. Parcijalnom integracijom $u=x$, $dv=\sin x\,dx$: $\displaystyle\int_{0}^{\pi/2} x\sin x\,dx=-x\cos x\Big|_{0}^{\pi/2}+\sin x\Big|_{0}^{\pi/2}=1$.

  \textbf{Rešenje:} a) $\dfrac{e^2+1}{4}$; b) $\dfrac{2\sqrt2-1}{3}$; c) $\dfrac12$ i $1$.
\end{enumerate}
\clearpage
\section*{Slot 10. Kombinatorika — rešenja}
\begin{enumerate}[label=\arabic*.]
\item \textbf{a)} Svaka od $6$ različitih kuglica se nezavisno šalje u jednu od $3$ kutije (pravilo proizvoda): $3^6=729$.

\textbf{b)} Potrebne su \emph{surjekcije} iz $6$ kuglica u $3$ kutije. Po formuli uključenja-isključenja
\[
\sum_{k=0}^{3}(-1)^k\binom{3}{k}(3-k)^6=3^6-3\cdot 2^6+3\cdot 1^6=729-192+3=540.
\]

\textbf{v)} Kuglice su jednake, kutije različite — „zvezde i pregrade”: broj nenegativnih rešenja $x_1+x_2+x_3=6$ jednak je $\binom{6+3-1}{3-1}=\binom{8}{2}=28$.

\textbf{g)} Kuglice su različite, kutije jednake, svaka neprazna — to je Stirlingov broj 2. vrste $S(6,3)$. Pri različitim kutijama rasporeda ima $540$, a permutacije $3$ jednakih kutija daju ponavljanje $3!$ puta, pa je
\[
S(6,3)=\frac{540}{3!}=\frac{540}{6}=90.
\]
Ako su pak i kuglice jednake (kutije jednake, svaka neprazna), potreban je broj razbijanja broja $6$ tačno na $3$ pozitivna sabirka: $6=4+1+1=3+2+1=2+2+2$, to jest $3$ načina.

\emph{Rešenje:} a) $729$; \; b) $540$; \; v) $28$; \; g) $90$ i $3$.

\item Broj nenegativnih rešenja jednačine $y_1+\cdots+y_4=N$ jednak je $\binom{N+3}{3}$.

\textbf{a)} Odmah $\binom{20+3}{3}=\binom{23}{3}=\dfrac{23\cdot 22\cdot 21}{6}=1771$.

\textbf{b)} Smena $x_i=y_i+1$, $y_i\ge 0$, daje $y_1+\cdots+y_4=16$, odakle $\binom{19}{3}=\dfrac{19\cdot 18\cdot 17}{6}=969$.

\textbf{v)} Smena $x_i=y_i+2$, $y_i\ge 0$, daje $y_1+\cdots+y_4=12$, odakle $\binom{15}{3}=\dfrac{15\cdot 14\cdot 13}{6}=455$.

\emph{Rešenje:} a) $1771$; \; b) $969$; \; v) $455$.

\item \textbf{a)} Prva cifra nije jednaka $0$: $5$ mogućnosti. Ostale tri pozicije — različite cifre iz preostalih: $5\cdot 4\cdot 3$. Ukupno $5\cdot 5\cdot 4\cdot 3=300$.

\textbf{b)} Deljivost sa $5$ znači da je poslednja cifra $0$ ili $5$. Ako je poslednja cifra $0$: prve tri pozicije — $5\cdot 4\cdot 3=60$. Ako je poslednja cifra $5$: prva nije $0$ ni $5$ ($4$ mogućnosti), zatim $4\cdot 3$, to jest $4\cdot 4\cdot 3=48$. Ukupno $60+48=108$.

\textbf{v)} Parnost: poslednja cifra $\in\{0,2,4\}$. Ako je poslednja $0$: $5\cdot 4\cdot 3=60$. Ako je poslednja $2$ ili $4$ ($2$ mogućnosti): prva nije $0$ ni izabrana poslednja ($4$ mogućnosti), dalje $4\cdot 3$, to jest $4\cdot 4\cdot 3=48$ za svaku, ukupno $96$. Ukupno $60+96=156$.

\emph{Rešenje:} a) $300$; \; b) $108$; \; v) $156$.

\item \textbf{a)} Rasporedimo za okruglim stolom $4$ mladića: kružnih permutacija $(4-1)!=3!=6$. Između mladića ima tačno $4$ razmaka, u koje sedimo $4$ devojke — $4!=24$ načina. Ukupno $6\cdot 24=144$.

\textbf{b)} Metoda razmaka. Rasporedimo u red ostalih $5$ osoba: $5!=120$ načina. One stvaraju $6$ razmaka (uključujući krajeve). U $3$ različita razmaka sedimo troje „nesuseda” uz uvažavanje redosleda: $6\cdot 5\cdot 4=120$. Ukupno $120\cdot 120=14400$.

\emph{Rešenje:} a) $144$; \; b) $14400$.

\item U reči ima $10$ slova: $A$ — tri puta, $M$ — dva puta, $T$ — dva puta, a $E,I,K$ — po jednom.

\textbf{a)} Permutacije sa ponavljanjem:
\[
\frac{10!}{3!\,2!\,2!}=\frac{3628800}{24}=151200.
\]

\textbf{b)} Slepimo tri $A$ u blok $[AAA]$ i permutujemo $8$ objekata (blok, $M,M,T,T,E,I,K$):
\[
\frac{8!}{2!\,2!}=\frac{40320}{4}=10080.
\]

\textbf{v)} Metoda razmaka. Rasporedimo $8$ slova bez dva $T$ ($A,A,A,M,M,E,I,K$): $\dfrac{8!}{3!\,2!}=3360$ načina. Ona stvaraju $9$ razmaka; u $2$ od njih stavljamo dva jednaka $T$: $\binom{9}{2}=36$. Ukupno $3360\cdot 36=120960$.

\emph{Rešenje:} a) $151200$; \; b) $10080$; \; v) $120960$.

\item \textbf{a)} Preko komplementa. Ukupno komisija od $4$ osobe: $\binom{12}{4}=495$. Komisija bez devojaka: $\binom{7}{4}=35$. „Bar jedna devojka”: $495-35=460$.

\textbf{b)} \emph{Kompozicije} (redosled je važan) broja $6$ na $3$ prirodna sabirka — broj pozitivnih rešenja $x_1+x_2+x_3=6$: $\binom{6-1}{3-1}=\binom{5}{2}=10$. \emph{Razbijanja} (redosled nije važan): $6=4+1+1=3+2+1=2+2+2$, to jest $3$.

\textbf{v)} Dva različita slova iz $6$ uz uvažavanje redosleda: $6\cdot 5=30$. Tri cifre, prva iz $\{1,\dots,9\}$, ostale dve bilo koje: $9\cdot 10\cdot 10=900$. Ukupno $30\cdot 900=27000$.

\emph{Rešenje:} a) $460$; \; b) kompozicija $10$, razbijanja $3$; \; v) $27000$.
\end{enumerate}

\end{document}
